% Equality-interface proof. Inclusion fragment.
% Requires amsmath, amssymb, amsthm and theorem/lemma/proposition environments.
% All labels use the prefix seed:.

\section{A deterministic equality rule with no template parameter}
\label{seed:main}

This reduction uses an initialized scalar iterator with its zero equation
made definitional. It differs from successor-seeded canonicalization:
here a zero-counter iterator disappears entirely. That feature lets a
single fixed-position equality rule include root equality composition.

\subsection{Canonical terms and the six rules}

Canonical terms have the constructors $x$, $0$, $S(t)$ and
$I(n,a;u.h)$, with canonical subterms and with the counter $n$ not
literally $0$. Only $u$ in $h$ is bound. Define the smart constructor
\[
 \mathcal I(n,a;u.h)=
 \begin{cases}
 a,&n\equiv0,\\
 I(n,a;u.h),&n\not\equiv0.
 \end{cases}
\]
The symbol $\equiv$ denotes literal syntactic identity. In this display
the arguments are already canonical. For arbitrary raw terms, normalize
all three arguments recursively before applying this smart constructor.
Thus normalization must test the \emph{normalized} counter.
Write $t\langle\sigma\rangle$ for capture-free substitution followed by
normalization. This operation is used throughout the calculus.

Native formulas are $s=t$, $Q_y B$, and $D_{x,y}B$. Their ordinary
readings are unchanged:
\[
 Q_yB=B(0)\to\forall y\,B(y),\qquad
 D_{x,y}B=\exists x\bigl(B(x,0)\to\forall y\,B(x,y)\bigr).
\]
All displayed instances of a matrix use canonical substitution.
Contexts are finite reusable lists of native formulas. Put $1=S0$.

For canonical terms $s,t,b$, define the partial operation
\[
 \rho_{s,t}(b)=
 \begin{cases}
 t,&b\equiv s,\\
 I(n,t;u.h),&b\equiv I(n,s;u.h),\\
 \text{undefined},&\text{otherwise}.
 \end{cases}
 \tag{Seed replacement}\label{seed:rho}
\]
The second branch inspects one literal constructor and its initializer;
it leaves the counter and update unchanged. The first and second branches
are disjoint. If both matched, the finite term $s$ would equal a term
having $s$ as a proper initializer subterm, which contradicts strict
increase of syntax-tree size. The output of the second branch remains
canonical because its counter is the same nonzero literal counter.
No bound variable is substituted or renamed by this operation.

Define
\[
 \eta_{B,t}(z)=
 \begin{cases}
 z,&B\text{ has equation root and }t\equiv0,\\
 0,&\text{otherwise}.
 \end{cases}
\]
The complete six-rule presentation, denoted $\mathcal S$, is:
\begin{align*}
 &\frac{A\in\Gamma}{\Gamma\vdash A}
 &&\text{Hyp},\\[3pt]
 &\frac{}{\Gamma\vdash
 I(Sn,a;u.h)=h\langle\mathcal I(n,a;u.h)/u\rangle}
 &&\text{IS},\\[3pt]
 &\frac{\Gamma\vdash s=t\qquad\Gamma\vdash b=r}
 {\Gamma\vdash \rho_{s,t}(b)=r}
 \quad(\rho_{s,t}(b)\text{ defined})
 &&\text{Seed},\\[3pt]
 &\frac{\Gamma\vdash Q_yB(y)\qquad\Gamma\vdash B(0)}
 {\Gamma\vdash B(t)}
 &&\text{QE},\\[3pt]
 &\frac{\Gamma\vdash D_{x,y}B(x,y)\qquad
 \Gamma,Q_yB(b,y)\vdash A}{\Gamma\vdash A}
 &&\text{DE},\\[3pt]
 &\frac{\Gamma,B(t,\eta_{B,t}(z))\vdash B(t,Sz)}
 {\Gamma\vdash D_{x,y}B(x,y)}
 &&\text{DI}.
\end{align*}
In DE, $b$ is fresh for the retained context, conclusion, and passive
matrix parameters. In DI, $z$ is fresh for the retained context, witness
$t$, and passive matrix parameters. Local assumptions are discharged.
The designated contradiction is $0=S0$.

The seed rule has no numerical side parameter in a proof certificate:
its two premise conclusions determine the conclusion, or determine
rejection. In particular it has neither an arbitrary numerical template
nor a user-supplied occurrence path. This does \emph{not} remove
substitution from the whole proof checker: IS and the binder rules still
require canonical, capture-free substitution.

\begin{lemma}[Equivalent displayed seed schema]\label{seed:schema}
The deterministic Seed rule is exactly the schema
\[
 \frac{\Gamma\vdash s=t\qquad
 \Gamma\vdash\mathcal I(n,s;u.h)=r}
 {\Gamma\vdash \mathcal I(n,t;u.h)=r},
 \tag{K}\label{seed:K}
\]
where $n,h$ are arbitrary canonical data. They need not be supplied in
certificates.
\end{lemma}
\begin{proof}
At $n=0$, the schema is the first branch of
\eqref{seed:rho}; the erased update is irrelevant, and may be chosen as
$0$. At any other canonical counter, the second premise displays
exactly the iterator node from which its counter and update are read.
The branches cannot overlap by the finite-size argument above. Conversely,
the first deterministic branch is the schema with $n=0,h=0$, and the
second is its instance with the counter and update displayed by $b$.
\end{proof}

\begin{lemma}[Seed replacement commutes with canonical substitution]
\label{seed:substitution}
If $\rho_{s,t}(b)$ is defined, then for every admissible numerical
substitution $\sigma$,
\[
 \rho_{s\langle\sigma\rangle,t\langle\sigma\rangle}
       (b\langle\sigma\rangle)
 =\rho_{s,t}(b)\langle\sigma\rangle.
\]
In particular the substituted partial operation is defined.
\end{lemma}
\begin{proof}
In the root branch this is literal equality. In the iterator branch,
canonical substitution changes $I(n,s;u.h)$ to the smart constructor
with substituted counter, initializer and update. If the substituted
counter normalizes to $0$, both the old minor and old conclusion lose
their iterator nodes, and the new rule uses its root branch. Otherwise
both retain their iterator nodes, and the new rule uses its initializer
branch. The counter and update are unchanged between the compared
expressions in either case. Canonical substitution coherence supplies
the stated identity. Thus the equality rule itself is stable even when
normalization changes its operational branch.
\end{proof}

\subsection{Complete equality reasoning from the seed rule}

Every argument in this subsection works under an arbitrary retained
context, including arbitrary equation hypotheses. It uses only IS and
Seed, apart from using the supplied premise derivations. It uses no
induction, numerical cases, or logical binder rule.

The instance of \eqref{seed:K} at $n=0$ is
\[
 \frac{s=t\qquad s=r}{t=r}.
 \tag{K0}\label{seed:K0}
\]
It is available for \emph{all} source terms, including literal $0$.
No source restriction is part of this presentation.

\begin{lemma}[Root equality algebra]\label{seed:algebra}
Reflexivity, symmetry and transitivity are derivable.
\end{lemma}
\begin{proof}
Put $L(a)=I(1,a;u.u)$. IS at predecessor zero gives $L(a)=a$,
because the right-hand substituted term is $\mathcal I(0,a;u.u)=a$.
Using two copies of this axiom as the premises of K0 gives $a=a$.

For symmetry of a supplied $s=t$, use it as K0's major and $s=s$
as its minor; the result is $t=s$. For transitivity of $s=t$ and
$t=r$, first obtain $t=s$ by symmetry. K0 with major $t=s$ and
minor $t=r$ gives $s=r$.
\end{proof}

\begin{lemma}[One-step application]\label{seed:beta}
For every fixed canonical numerical template $f(u)$ and ambient term $a$,
IS derives
\[
 I(1,a;u.f(u))=f\langle a/u\rangle.
 \tag{Beta1}\label{seed:betaeq}
\]
\end{lemma}
\begin{proof}
IS with predecessor zero has right side
$f\langle\mathcal I(0,a;u.f)/u\rangle$, which is literally
$f\langle a/u\rangle$. No equality substitution into $f$ has been
used. This is the point at which definitional zero computation is essential.
\end{proof}

\begin{theorem}[Arbitrary numerical congruence]\label{seed:congruence}
For every fixed numerical template $f(u)$, the rule
\[
 \frac{\Gamma\vdash s=t}
 {\Gamma\vdash f\langle s/u\rangle=f\langle t/u\rangle}
\]
is admissible in $\mathcal S$.
\end{theorem}
\begin{proof}
Abbreviate $f\langle s/u\rangle$ by $f(s)$ and similarly for $t$,
and put $A_a=I(1,a;u.f)$.
The entire construction is
\[
 \frac{
   \displaystyle\frac{s=t\qquad A_s=f(s)}{A_t=f(s)}\;\mathrm{Seed}_1
   \qquad A_t=f(t)
 }{f(s)=f(t)}\;\mathrm{Seed}_0.
\]
The two endpoint equations $A_s=f(s)$ and $A_t=f(t)$ are IS axioms
from Lemma~\ref{seed:beta}. The last inference is K0.
In particular this construction does not even use the derived
symmetry, transitivity, or reflexivity lemmas.
The template may contain zero, one, or several occurrences of its
designated variable, including occurrences under nested iterator binders.
All such occurrences are handled by IS's ordinary capture-free
substitution, not by the primitive equality rule.
\end{proof}

\begin{proposition}[Four-node congruence expansion]\label{seed:nodes}
If the supplied derivation of $s=t$ has $P$ inference nodes, the displayed
derivation of $f(s)=f(t)$ has exactly $P+4$ nodes: two IS axioms and two
Seed applications are added, and the supplied premise is inserted once.
The bound is independent of the syntactic depth and number of occurrences
of the distinguished variable in $f$.
\end{proposition}
\begin{proof}
Count the four displayed additions in the preceding proof tree.
\end{proof}

This is a count of the specified inference tree. Its term labels contain
copies of $f$, $s$, and $t$; it is not a constant bound on serialized
certificate size, runtime, or the shortest possible proof. The prior
wrapped counter-congruence rule in
\texttt{witness\_rule/theory/restricted\_congruence.tex} has a specified
$P+5$ construction using one counter-congruence inference, two endpoint
computations, and two separate root replacements. The present rule fuses
the lifting and the first root replacement at the inference interface;
the zero-counter instance of that same rule supplies root replacement.
The numerical language and its normalization convention differ, so this
comparison should not be read as a uniform shortening of every old proof.

\begin{corollary}[One-template transport]\label{seed:onesided}
The admissible rule
\[
 \frac{s=t\qquad f(s)=r}{f(t)=r}
\]
holds with arbitrary source terms and arbitrary contexts.
\end{corollary}
\begin{proof}
Theorem~\ref{seed:congruence} supplies $f(s)=f(t)$. K0 with this
major and the supplied minor gives the result.
\end{proof}

\begin{corollary}[One-template crossed transport]\label{seed:crossed}
The admissible rule
\[
 \frac{s=t\qquad f(s)=r}{r=f(t)}
\]
holds with arbitrary source terms and arbitrary contexts.
\end{corollary}
\begin{proof}
Corollary~\ref{seed:onesided} gives $f(t)=r$; apply derived symmetry.
\end{proof}

\begin{corollary}[Full equational transport]\label{seed:eqtransport}
For arbitrary numerical templates $a(u),b(u)$,
\[
 \frac{s=t\qquad a(s)=b(s)}{a(t)=b(t)}
\]
is admissible.
\end{corollary}
\begin{proof}
Apply Corollary~\ref{seed:crossed} first with $f=a,r=b(s)$,
obtaining $b(s)=a(t)$, and then with $f=b,r=a(t)$, obtaining
$a(t)=b(t)$. These are two uses of an admissible rule, not two
primitive seed instructions; their numerical-congruence expansions
contain computation and root-composition steps.
\end{proof}

Weakening, fresh renaming, and ordinary-assumption cut in $\mathcal S$
follow directly by induction on derivations. They leave all numerical
expressions at a rule instance unchanged, up to injective renaming.
Thus neither the DI mode nor the deterministic Seed branch changes.
General numerical substitution is addressed after the comparison theorem;
it should not be claimed by a direct induction on restricted DI.

\subsection{Recovering both ordinary introductions}

\begin{lemma}[A canonical nonliteral zero]\label{seed:alias}
The term
\[
 Z:=I(1,0;u.0)
\]
is canonical, is not literally $0$, and has an IS proof of $Z=0$.
\end{lemma}
\begin{proof}
Its counter is $S0$, so canonicalization does not erase its outer node.
IS with predecessor zero and constant-zero update has right side $0$.
\end{proof}

The old zero-counter alias $I(0,0;u.0)$ cannot serve here: it is literally
$0$ in the canonical syntax. The one-step alias in
Lemma~\ref{seed:alias} is the required replacement.

\begin{lemma}[Unrestricted case introduction]\label{seed:case}
For every matrix $B$, every witness $t$, and the usual fresh $z$,
\[
 \frac{\Gamma,B(t,0)\vdash B(t,Sz)}
 {\Gamma\vdash D_{x,y}B(x,y)}
 \tag{D-case}\label{seed:Dcase}
\]
is admissible in $\mathcal S$.
\end{lemma}
\begin{proof}
The displayed instance already is primitive DI unless $B$ is an equation
and $t$ is literally zero. In the exceptional case write $B=E$ and
let the supplied derivation be
\[
 \pi:\quad\Gamma,E(0,0)\vdash E(0,Sz).
\]
Choose $Z$ from Lemma~\ref{seed:alias}. Under
$\Gamma,E(Z,0)$, transport its local hypothesis along $Z=0$ using
Corollary~\ref{seed:eqtransport}, obtaining $E(0,0)$.
Ordinary-assumption cut with $\pi$ gives $E(0,Sz)$.
Reverse-transport along $0=Z$ to obtain $E(Z,Sz)$.
DI at the nonliteral witness $Z$ is in its case mode and concludes
the desired $D_{x,y}E(x,y)$.

Only equation transport is used, since $E$ has equation root.
No transport through a binder, numerical substitution into $\pi$, or
substitution-admissibility theorem for the restricted calculus is assumed.
\end{proof}

\begin{lemma}[Direct equation progress]\label{seed:progress}
For every literal equation $E(y)$,
\[
 \frac{\Gamma,E(z)\vdash E(Sz)}{\Gamma\vdash Q_yE(y)}
 \tag{Q-progress}\label{seed:Qprogress}
\]
is admissible in $\mathcal S$.
\end{lemma}
\begin{proof}
Choose a witness variable $x$ absent from $E$ and apply primitive DI
to $D_{x,y}E(y)$ at witness $0$. Its mode is progress, and its premise
is precisely the given one. DE opens this formula with a fresh witness;
the resulting local assumption is $Q_yE(y)$, independent of that witness.
Hyp proves it in the opened context and DE exports it.
\end{proof}

\begin{lemma}[Conditional closure introduction]\label{seed:Qcase}
For arbitrary $B$, D-case and DE derive
\[
 \frac{\Gamma,B(0)\vdash B(Sz)}{\Gamma\vdash Q_yB(y)}.
 \tag{Q-case}\label{seed:qcase}
\]
\end{lemma}
\begin{proof}
Use D-case at an unused witness coordinate and then DE plus Hyp,
exactly as in Lemma~\ref{seed:progress}.
\end{proof}

\subsection{The stable auxiliary calculus and formula transport}

Let $\mathcal A$ use the same canonical terms, formulas, and contexts.
Its rules are Hyp, IS, unrestricted equational transport,
QE, DE, D-case, and Q-progress. This is the canonical counterpart of
the native kernel with zero computation made definitional.

\begin{theorem}[Exact sequent equivalence]\label{seed:auxequiv}
For every finite context $\Gamma$ and native formula $A$,
\[
 \Gamma\vdash_{\mathcal S}A
 \quad\Longleftrightarrow\quad
 \Gamma\vdash_{\mathcal A}A.
\]
\end{theorem}
\begin{proof}
From $\mathcal A$ to $\mathcal S$, recursively translate premise
derivations. Corollary~\ref{seed:eqtransport} replaces equational
transport, Lemma~\ref{seed:case} replaces D-case, and
Lemma~\ref{seed:progress} replaces Q-progress. Other rules are unchanged.

Conversely, the standard equality algebra is derivable in
$\mathcal A$ from IS and unrestricted equational transport, starting
again with the computation $L(a)=a$. An instance of K follows by
transport along $s=t$ inside the single numerical template
$\mathcal I(n,v;u.h)$, keeping $n,h$ fixed.
This also covers $n=0$, when that canonical template is just $v$.
Each DI case-mode inference is already D-case.

For a progress-mode DI, the witness is $0$ and its matrix is an equation
$E(x,y)$. Q-progress applied to the supplied premise gives
$Q_yE(0,y)$. Under $E(0,0)$, QE at $Sz$ yields $E(0,Sz)$.
D-case with witness $0$ now concludes $D_{x,y}E(x,y)$.
This translates every remaining primitive rule and preserves end sequents.
\end{proof}

\begin{theorem}[Full formula transport]\label{seed:fulltransport}
For every native formula template $F(v)$,
\[
 s=t,\ F(s)\ \vdash_{\mathcal S}\ F(t).
\]
\end{theorem}
\begin{proof}
Use external structural induction on $F$, uniformly in the retained
context and all passive parameters. Equation roots use
Corollary~\ref{seed:eqtransport}.

Suppose $F(v)=Q_yB(v,y)$. To introduce $Q_yB(t,y)$ by Q-case,
assume $B(t,0)$. Symmetry and the induction hypothesis convert this
to $B(s,0)$. QE for the given $Q_yB(s,y)$ gives $B(s,Sz)$,
and the induction hypothesis transports this to $B(t,Sz)$.
Discharge the base using Q-case.

Suppose $F(v)=D_{x,y}B(v,x,y)$. DE opens the source at a fresh
witness $b$, supplying $Q_yB(s,b,y)$. Introduce the target by
D-case at the same witness. Its local base $B(t,b,0)$ is transported
backwards to $B(s,b,0)$, QE gives $B(s,b,Sz)$, and forward transport
gives $B(t,b,Sz)$. Close D-case and DE. Every recursive call concerns
the strictly smaller matrix $B$; no arithmetic induction on a compound
formula has been introduced.
\end{proof}

The auxiliary calculus is stable under canonical numerical substitution:
normalization coherence handles IS and all instances, formula roots are
preserved, and eigenvariables are freshened. Theorem~\ref{seed:auxequiv}
therefore gives substitution admissibility for $\mathcal S$ as well.
This is a derived admissibility result. A substitution can change a
witness into literal $0$, so a direct rule-for-rule substitution proof
for its six-rule presentation would be invalid.

\subsection{Interpretations and the precise novelty claim}

Once the raw/canonical comparison is established, the existing native
kernel's whole-domain bi-interpretation with iPRA transfers by
Theorem~\ref{seed:auxequiv}. The numerical domains remain whole,
equality remains actual equality, and the numerical composite
isomorphisms have graph $x=y$. The raw/canonical comparison must prove
normalization correctness under iterator bodies; it cannot replace that
obligation with a purely syntactic claim that the ordinary theories are
identical. Literal equation induction supplies the needed body
extensionality. Section~\ref{zeroI:section} proves that comparison and its full formula composites.

The existing sources already contain several ingredients:
\begin{itemize}
\item \texttt{Research/counter\_congruence/theory/counter\_theorem.tex}
derives arbitrary context congruence from root replacement plus
counter-position congruence for a \emph{counter-aware}, successor-seeded
recursor.
\item \texttt{Research/witness\_rule/theory/restricted\_congruence.tex}
further gives a wrapped, state-independent counter rule, a separate root
rule, and a five-node congruence expansion. It explicitly factors arbitrary
templates through a fixed computation shape.
\item \texttt{Research/one\_context} and the incoming
\texttt{ipra\_one\_sided\_transport\_2026-11-11.zip} use one-template
crossed transport. The single witness-directed introduction is inherited
from \texttt{witness\_rule} and \texttt{NativeOneIntro.lean}.
\item The incoming \texttt{ipra\_canonical\_recursors\_2026-11-12.zip}
makes the zero computation of a successor-seeded recursor canonical by
erasing its update body. It does not erase the whole recursor node and
does not make its zero step the identity function of the initializer.
\end{itemize}

The present specific change is that \emph{one initialized scalar seed
rule}, with zero computation made definitional, supplies both root
composition (counter $0$) and arbitrary context congruence (counter $1$).
Its conclusion is determined by its two premise conclusions. No template,
term, or occurrence path need be supplied to the equality instruction.
This is a mathematical reduction of the equality interface, followed by
an implementation opportunity. It is not a claim that fixed-position
congruence is a new idea, that normalization has no trusted cost, or that
an absolute smallest iPRA axiomatization has been found.

All source paths above refer to the commit stated in
Section~\ref{sec:assessment}, under
\path{Logic/PrimitiveRecursiveArithmetic/ConciseIPRA}.
